\documentclass[12pt]{article}
\usepackage{amsmath}
\usepackage{amssymb}
\usepackage{graphicx}
\usepackage{hyperref}
\usepackage{geometry}
\geometry{margin=1in}

\title{Pintle Injector Liquid Rocket Engine:\\
Physics, Mathematics, and Dynamics}
\author{Kush Mahajan}
\date{\today}

\begin{document}

\maketitle

\tableofcontents
\newpage

\section{Introduction}

This document provides a comprehensive mathematical and physical description of the pintle injector liquid rocket engine simulation pipeline. The pipeline takes tank pressures as input and solves for chamber pressure, mass flow rates, and thrust output through a coupled system of equations.

\subsection{Workflow Overview}

The pipeline follows this sequence:
\begin{enumerate}
    \item \textbf{Input}: Tank pressures $P_{\text{tank},O}$ and $P_{\text{tank},F}$
    \item \textbf{Feed System}: Calculate pressure losses from tank to injector
    \item \textbf{Injector Flow}: Calculate mass flow rates through pintle injector
    \item \textbf{Spray Physics}: Validate atomization and mixing quality
    \item \textbf{Chamber Solver}: Find equilibrium chamber pressure $P_c$
    \item \textbf{Nozzle}: Calculate thrust from exit conditions
\end{enumerate}

\section{Pintle Injector Geometry}

\subsection{LOX Injector (Axial Flow)}

LOX flows through $N$ orifices on the pintle tip. Each orifice has diameter $d_{\text{orifice}}$.

\textbf{Total LOX Flow Area:}
\begin{equation}
    A_{\text{LOX}} = N \times \frac{\pi d_{\text{orifice}}^2}{4}
\end{equation}

\textbf{Hydraulic Diameter:}
\begin{equation}
    d_{h,\text{LOX}}} = d_{\text{orifice}}
\end{equation}

\subsection{Fuel Injector (Radial Flow)}

Fuel flows through an annulus gap between the pintle tip and reservoir wall. The pintle tip has outer diameter $d_{\text{pintle}}$, and the gap height is $h_{\text{gap}}$.

\textbf{Inner Radius:}
\begin{equation}
    R_{\text{inner}} = \frac{d_{\text{pintle}}}{2}
\end{equation}

\textbf{Outer Radius:}
\begin{equation}
    R_{\text{outer}} = R_{\text{inner}} + h_{\text{gap}}
\end{equation}

\textbf{Fuel Flow Area (Annulus):}
\begin{equation}
    A_{\text{fuel}} = \pi \left( R_{\text{outer}}^2 - R_{\text{inner}}^2 \right)
\end{equation}

\textbf{Hydraulic Diameter (Thin Annulus):}
\begin{equation}
    d_{h,\text{fuel}} = 2 h_{\text{gap}}
\end{equation}

\section{Feed System Pressure Losses}

Pressure losses in the feed system from tank to injector are modeled using a generalized loss coefficient:

\begin{equation}
    \Delta p_{\text{feed}} = K_{\text{eff}}(P) \times \frac{\rho}{2} \times \left( \frac{\dot{m}}{\rho A_{\text{hyd}}} \right)^2
\end{equation}

where:
\begin{itemize}
    \item $\dot{m}$ is the mass flow rate [kg/s]
    \item $\rho$ is the fluid density [kg/m³]
    \item $A_{\text{hyd}}$ is the hydraulic area of the feed line [m²]
    \item $K_{\text{eff}}(P)$ is the effective loss coefficient
\end{itemize}

\subsection{Effective Loss Coefficient}

The effective loss coefficient depends on pressure:

\begin{equation}
    K_{\text{eff}}(P) = K_0 + K_1 \times \phi(P)
\end{equation}

where $\phi(P)$ can be:
\begin{itemize}
    \item \textbf{None}: $\phi(P) = 0$ (constant $K_{\text{eff}} = K_0$)
    \item \textbf{Square root}: $\phi(P) = \sqrt{P}$
    \item \textbf{Logarithmic}: $\phi(P) = \ln(P)$
\end{itemize}

\subsection{Regenerative Cooling (Fuel Only)}

For the fuel line, additional pressure drop occurs through regenerative cooling channels:

\begin{equation}
    \Delta p_{\text{regen}} = \Delta p_{\text{inlet}} + \Delta p_{\text{split}} + \Delta p_{\text{channels}} + \Delta p_{\text{merge}} + \Delta p_{\text{outlet}}
\end{equation}

\subsubsection{Inlet Pipe}

Friction factor using Colebrook-White (Haaland approximation for rough pipes):

\begin{equation}
    f = \left( -1.8 \log_{10} \left[ \left( \frac{\epsilon/D}{3.7} \right)^{1.11} + \frac{6.9}{\text{Re}} \right] \right)^{-2}
\end{equation}

For smooth pipes:
\begin{equation}
    f = \begin{cases}
        \frac{64}{\text{Re}} & \text{if Re} < 2300 \text{ (laminar)} \\
        \frac{0.316}{\text{Re}^{0.25}} & \text{if Re} \geq 2300 \text{ (turbulent, Blasius)}
    \end{cases}
\end{equation}

Pressure drop:
\begin{equation}
    \Delta p_{\text{inlet}} = f \frac{L}{D} \frac{\rho u^2}{2}
\end{equation}

\subsubsection{Parallel Channels}

For $N$ parallel channels, mass flow per channel:
\begin{equation}
    \dot{m}_{\text{channel}} = \frac{\dot{m}_{\text{total}}}{N}
\end{equation}

Rectangular channel hydraulic diameter:
\begin{equation}
    D_h = \frac{2wh}{w + h}
\end{equation}

where $w$ is width and $h$ is height.

\section{Injector Mass Flow}

Mass flow through the injector is calculated using the orifice flow equation:

\begin{equation}
    \dot{m} = C_d \times A \times \sqrt{2 \rho \Delta p_{\text{injector}}}
\end{equation}

where:
\begin{itemize}
    \item $C_d$ is the discharge coefficient
    \item $A$ is the flow area (either $A_{\text{LOX}}$ or $A_{\text{fuel}}$)
    \item $\rho$ is the fluid density
    \item $\Delta p_{\text{injector}} = P_{\text{injector}} - P_c$ is the pressure drop across the injector
\end{itemize}

\subsection{Discharge Coefficient}

The discharge coefficient depends on Reynolds number:

\begin{equation}
    C_d(\text{Re}) = C_{d,\infty} - \frac{a_{\text{Re}}}{\sqrt{\text{Re}}}
\end{equation}

with constraints:
\begin{equation}
    C_{d,\min} \leq C_d \leq C_{d,\infty}
\end{equation}

For fixed $C_d$ (matching spreadsheet approach), set $a_{\text{Re}} = 0$.

\subsection{Reynolds Number}

\begin{equation}
    \text{Re} = \frac{\rho u D_h}{\mu}
\end{equation}

where:
\begin{itemize}
    \item $u$ is the flow velocity [m/s]
    \item $D_h$ is the hydraulic diameter [m]
    \item $\mu$ is the dynamic viscosity [Pa·s]
\end{itemize}

Flow velocity:
\begin{equation}
    u = \frac{\dot{m}}{\rho A}
\end{equation}

\section{Spray Physics}

\subsection{Momentum Flux Ratio}

The momentum flux ratio $J$ characterizes the relative momentum of the two streams:

\begin{equation}
    J = \frac{\rho_O u_O^2}{\rho_F u_F^2}
\end{equation}

where subscripts $O$ and $F$ denote oxidizer and fuel, respectively.

\subsection{Spray Angle}

The spray angle $\theta$ can be calculated from either $J$ or the thrust momentum ratio (TMR):

\textbf{From J:}
\begin{equation}
    \theta = k J^n
\end{equation}

\textbf{From TMR:}
\begin{equation}
    \text{TMR} = \frac{J}{1 + \text{MR}}
\end{equation}

where MR is the mixture ratio (O/F).

\subsection{Weber Number}

The Weber number characterizes atomization quality:

\begin{equation}
    \text{We} = \frac{\rho u^2 d}{\sigma}
\end{equation}

where:
\begin{itemize}
    \item $d$ is the characteristic diameter (orifice or hydraulic)
    \item $\sigma$ is the surface tension [N/m]
\end{itemize}

Minimum Weber number constraint ensures proper atomization:
\begin{equation}
    \text{We} \geq \text{We}_{\min}
\end{equation}

\subsection{Sauter Mean Diameter (SMD)}

The Sauter Mean Diameter $D_{32}$ is calculated using Lefebvre's correlation:

\begin{equation}
    D_{32} = d \times C \times \text{We}^{-m} \times \text{Oh}^p
\end{equation}

where:
\begin{itemize}
    \item $C$, $m$, $p$ are empirical constants
    \item $\text{Oh} = \frac{\mu}{\sqrt{\rho \sigma d}}$ is the Ohnesorge number
\end{itemize}

\subsection{Evaporation Length}

The evaporation length $x^*$ estimates how far droplets travel before fully evaporating:

\begin{equation}
    x^* = U_{\text{rel}} \times \tau_{\text{evap}}
\end{equation}

where:
\begin{itemize}
    \item $U_{\text{rel}} = \sqrt{u_O^2 + u_F^2}$ is the relative velocity magnitude
    \item $\tau_{\text{evap}} = K \times D_{32}^2$ is the evaporation timescale
    \item $K$ is an empirical constant
\end{itemize}

Maximum evaporation length constraint ensures complete combustion:
\begin{equation}
    x^* \leq x^*_{\max}
\end{equation}

\section{Combustion Efficiency}

\subsection{Characteristic Length}

The characteristic length $L^*$ relates chamber volume to throat area:

\begin{equation}
    L^* = \frac{V_{\text{chamber}}}{A_{\text{throat}}}
\end{equation}

\subsection{Combustion Efficiency Model}

Combustion efficiency $\eta_{c^*}$ accounts for finite chamber effects (L*-based correction):

\begin{equation}
    \eta_{c^*} = 1 - C \times e^{-K \times L^*}
\end{equation}

where $C$ and $K$ are empirical coefficients (default from Huzel \& Huang).

\subsection{Actual Characteristic Velocity}

The actual characteristic velocity accounts for combustion efficiency:

\begin{equation}
    c^*_{\text{actual}} = \eta_{c^*} \times c^*_{\text{ideal}}
\end{equation}

where $c^*_{\text{ideal}}$ comes from Chemical Equilibrium Analysis (CEA).

\subsection{Actual Chamber Temperature}

The actual chamber temperature is estimated from ideal temperature and efficiency:

\begin{equation}
    T_{c,\text{actual}} = T_{c,\text{ideal}} \times \left( \frac{\eta_{c^*}}{1 - (1 - \eta_{c^*}) \times \frac{\gamma - 1}{\gamma}} \right)
\end{equation}

\subsection{Frozen Flow Correction}

A correction factor for $\gamma$ accounts for frozen flow effects in the nozzle:

\begin{equation}
    \gamma_{\text{actual}} = \gamma_{\text{ideal}} \times \left( 1 - \alpha \times (1 - \eta_{c^*}) \right)
\end{equation}

where $\alpha$ is an empirical constant.

\section{Chamber Pressure Solver}

The chamber pressure $P_c$ is solved by balancing propellant supply and combustion demand.

\subsection{Supply Side}

Mass flow supply from injectors:

\begin{equation}
    \dot{m}_{\text{supply}} = \dot{m}_O + \dot{m}_F
\end{equation}

where $\dot{m}_O$ and $\dot{m}_F$ are calculated from injector flow equations, which depend on $P_c$ through:
\begin{equation}
    \Delta p_{\text{injector}} = P_{\text{injector}} - P_c
\end{equation}

\subsection{Demand Side}

Mass flow demand from combustion:

\begin{equation}
    \dot{m}_{\text{demand}} = \frac{P_c \times A_{\text{throat}}}{c^*_{\text{actual}}}
\end{equation}

where $c^*_{\text{actual}}$ depends on $P_c$ and the mixture ratio MR through CEA.

\subsection{Residual Function}

The solver finds $P_c$ such that:

\begin{equation}
    R(P_c) = \dot{m}_{\text{supply}}(P_c) - \dot{m}_{\text{demand}}(P_c) = 0
\end{equation}

This is solved using Brent's method (root-finding algorithm).

\subsection{Bounds}

The search bounds are:
\begin{itemize}
    \item \textbf{Lower bound}: $P_{c,\min} = 0.1 \times \min(P_{\text{tank},O}, P_{\text{tank},F})$
    \item \textbf{Upper bound}: $P_{c,\max} = 0.95 \times \min(P_{\text{tank},O}, P_{\text{tank},F})$
\end{itemize}

\section{Nozzle Thrust Calculation}

\subsection{Thrust Components}

Total thrust consists of two components:

\begin{equation}
    F = F_{\text{momentum}} + F_{\text{pressure}}
\end{equation}

\subsubsection{Momentum Thrust}

\begin{equation}
    F_{\text{momentum}} = \dot{m}_{\text{total}} \times v_{\text{exit}}
\end{equation}

\subsubsection{Pressure Thrust}

\begin{equation}
    F_{\text{pressure}} = (P_{\text{exit}} - P_a) \times A_{\text{exit}}
\end{equation}

where $P_a$ is the ambient pressure.

\subsection{Exit Conditions}

Exit conditions are calculated using isentropic relations.

\subsubsection{Exit Mach Number}

From area ratio:
\begin{equation}
    \frac{A_{\text{exit}}}{A_{\text{throat}}} = \frac{1}{M_{\text{exit}}} \left( \frac{2}{\gamma + 1} \left( 1 + \frac{\gamma - 1}{2} M_{\text{exit}}^2 \right) \right)^{\frac{\gamma + 1}{2(\gamma - 1)}}
\end{equation}

This is solved iteratively for $M_{\text{exit}}$.

\subsubsection{Exit Pressure}

\begin{equation}
    \frac{P_{\text{exit}}}{P_c} = \left( 1 + \frac{\gamma - 1}{2} M_{\text{exit}}^2 \right)^{-\frac{\gamma}{\gamma - 1}}
\end{equation}

\subsubsection{Exit Temperature}

\begin{equation}
    \frac{T_{\text{exit}}}{T_c} = \left( 1 + \frac{\gamma - 1}{2} M_{\text{exit}}^2 \right)^{-1}
\end{equation}

\subsubsection{Exit Velocity}

\begin{equation}
    v_{\text{exit}} = M_{\text{exit}} \times \sqrt{\gamma R T_{\text{exit}}}
\end{equation}

where $R$ is the specific gas constant.

\subsection{Specific Impulse}

\begin{equation}
    I_{\text{sp}} = \frac{F}{\dot{m}_{\text{total}} \times g_0}
\end{equation}

where $g_0 = 9.80665$ m/s² is standard gravity.

\section{Iterative Closure Algorithm}

The mass flow rates and chamber pressure are solved iteratively:

\begin{enumerate}
    \item \textbf{Initial guess}: $P_c^{(0)} = 0.5 \times \min(P_{\text{tank},O}, P_{\text{tank},F})$
    \item \textbf{For each iteration} $k$:
    \begin{enumerate}
        \item Calculate feed losses: $\Delta p_{\text{feed}}^{(k)}$
        \item Calculate injector inlet pressures: $P_{\text{injector}}^{(k)} = P_{\text{tank}} - \Delta p_{\text{feed}}^{(k)}$
        \item Calculate injector pressure drops: $\Delta p_{\text{injector}}^{(k)} = P_{\text{injector}}^{(k)} - P_c^{(k)}$
        \item Calculate mass flow rates: $\dot{m}_O^{(k)}$, $\dot{m}_F^{(k)}$
        \item Calculate Reynolds numbers and update $C_d$
        \item Check spray constraints (We, $x^*$)
        \item If constraints violated, reduce effective $C_d$ and recalculate
        \item Calculate mixture ratio: $\text{MR}^{(k)} = \dot{m}_O^{(k)} / \dot{m}_F^{(k)}$
        \item Get CEA properties: $c^*_{\text{ideal}}^{(k)}$, $\gamma^{(k)}$, etc.
        \item Calculate combustion efficiency: $\eta_{c^*}^{(k)}$
        \item Calculate actual $c^*$: $c^*_{\text{actual}}^{(k)} = \eta_{c^*}^{(k)} \times c^*_{\text{ideal}}^{(k)}$
        \item Calculate demand: $\dot{m}_{\text{demand}}^{(k)} = P_c^{(k)} \times A_{\text{throat}} / c^*_{\text{actual}}^{(k)}$
        \item Calculate residual: $R^{(k)} = \dot{m}_{\text{supply}}^{(k)} - \dot{m}_{\text{demand}}^{(k)}$
        \item Update $P_c^{(k+1)}$ using root-finding algorithm
    \end{enumerate}
    \item \textbf{Convergence}: When $|R^{(k)}| < \epsilon$ (tolerance)
\end{enumerate}

\section{Chemical Equilibrium Analysis (CEA)}

CEA provides thermochemical properties as a function of mixture ratio MR and chamber pressure $P_c$:

\begin{itemize}
    \item $c^*_{\text{ideal}}$: Ideal characteristic velocity [m/s]
    \item $C_{f,\text{ideal}}$: Ideal thrust coefficient
    \item $T_c$: Chamber temperature [K]
    \item $\gamma$: Ratio of specific heats
    \item $R$: Specific gas constant [J/(kg·K)]
    \item $M$: Molecular weight [kg/kmol]
\end{itemize}

These are pre-computed in a lookup table and interpolated during simulation for performance.

\section{Summary}

The pintle engine pipeline solves a coupled system where:
\begin{itemize}
    \item Mass flow rates depend on chamber pressure (through injector pressure drop)
    \item Chamber pressure depends on mass flow rates (through supply/demand balance)
    \item The system is solved iteratively using root-finding algorithms
    \item All physics (feed losses, injector flow, spray, combustion, nozzle) are coupled
\end{itemize}

The key insight is that \textbf{chamber pressure is solved, not input}. Given tank pressures, the system finds the equilibrium operating point where propellant supply matches combustion demand.

\end{document}

